Un moteur fait tourner une tige homogène de longueur $L=AB=\qty{0,5}{\meter}$
autour d'un axe fixe ($\Delta$) passant par son centre $O$
(figure~\ref{fig:tige}). Un point $M$ appartient à la tige est distant de l'axe
de rotation par la distance $OM=\frac{AB}{4}$. La
figure~\ref{fig:tige_theta_fct_t} représente la variation de l'abscisse
angulaire du point $M$ en fonction du temps.
\begin{figure}[H]
    \centering
    \begin{subfigure}{0.49\textwidth}
        \centering
        \begin{tikzpicture}[label distance=1pt]
            \newdimen\L \L=6cm
            \newdimen\OM \OM=\L
            \divide\OM by 4
            \def\a{45}
            \node[draw,fill=gray,inner sep=0pt,minimum width=\L,
                minimum height=4mm,rotate=\a] (T) {};
            \node[circle,densely dashed,draw,inner sep=0pt,
                minimum size=2\OM] (C) {};
            \foreach \l/\p in {O/(T.center),M/(C.\a),A/(T.east),B/(T.west)}
            \node[circle,draw,inner sep=0.8pt,label=below right:$\l$] at \p {};
            \draw[-Stealth] (90:1.4\OM) arc (90:90+30:1.4\OM)
                node[above,pos=0.5] {$\oplus$};
            \node[above left=1mm] {$\Delta$};
        \end{tikzpicture}
        \caption{}
        \label{fig:tige}
    \end{subfigure}
    \hspace{1mm}
    \begin{subfigure}{0.49\textwidth}
        \centering
        %\begin{tikzpicture}
        %    \draw[help lines] (0,0) grid (4.5,5.5);
        %    \draw[-Latex] (0,0) -- (4.5,0) node[right] {$t~(\unit{\ms})$};
        %    \draw[-Latex] (0,0) -- (0,5.5) node[right] {$\theta~(\unit{\rad})$};
        %    \foreach \x/\l in {0/0,1/10,2/20,3/30,4/40}
        %        \node[below] at (\x,0) {$\l$}; 
        %    \foreach \y/\l in {1/\pi,2/2\pi,3/3\pi,4/4\pi,5/5\pi}
        %        \node[left] at (0,\y) {$\l$};
        %    \draw[thick] (0,1) -- (4,5);
        %\end{tikzpicture}
        \begin{tikzpicture}[]
            \begin{axis}[
                    cycle list name=my color list,
                    width=7cm,height=7cm,
                    xlabel=$t~(\unit{\ms})$,ylabel=$\theta~(\unit{\rad})$,
                    xmin=0,xmax=55,
                    ymin=0,ymax=5.5*pi,
                    xtick distance=10,
                    ytick={0,pi,2*pi,3*pi,4*pi,5*pi},
                    yticklabels={$0$,$\pi$,$2\pi$,$3\pi$,$4\pi$,$5\pi$},
                    variable=t,
                ]
                \addplot[domain=0:40] {pi/10*t+pi};
            \end{axis}
        \end{tikzpicture}
        \caption{}
        \label{fig:tige_theta_fct_t}
    \end{subfigure}
    \caption{}
\end{figure}
\begin{enumerate}
    \item Quelle est la nature du mouvement de rotation de la tige ? Justifier
          votre réponse.
    \item Déterminer graphiquement la vitesse angulaire $\omega$ de la tige et
          la valeur de l'abscisse angulaire $\theta_{0}$ à $t=\qty{0}{\second}$.
    \item Écrire l'équation horaire $\theta(t)$ du mouvement de la tige.
    \item Calculer la période $T$ du mouvement de la tige, et déduire sa
          fréquence $f$.
    \item Calculer la durée $\Delta t$ nécessaire pour que la tige effectuée 20
          tours.
    \item Calculer la vitesse linéaire $v$ du point $M$ et la valeur de son
          abscisse curviligne $s_{0}$ à $t=\qty{0}{\second}$.
    \item Écrire l'équation horaire de l'abscisse curviligne $s(t)$ du point
          $M$.
\end{enumerate}

